\section{Conclusion}\label{sec:conclusion}

The bidder pool in a takeover depends on the information that buyers have before they commit to preparation. A
stronger incumbent moves the class value gap from the challenger's profit to target proceeds. In the order slot, a
liquidity order in the same direction is always possible. An investor who knows the challenger's class therefore
keeps an advantage after every order. A larger spread can make that advantage worth the cost of trade. The price then
recruits a challenger that a weaker incumbent could not attract.

The comparison holds with unique outcomes on an open set of primitives, and the admissible region is exact at the
benchmark. Between the two economies, informative and uninformative equilibria coexist. Along the informative branch,
entry rises weakly with incumbent strength, and strictly where the investor mixes. When information is costly,
competition creates research: research and entry rise together along a branch beside the equilibrium without
research.

Sale terms affect participation through the buyer's expected payment and through the information in the stock price.
In the recruitment region, access to prices raises target proceeds and acquisition surplus net of preparation costs.
The weight that the payment rule places on the runner-up's value decides how competition moves the spread. A reserve
that excludes the low class can make the price informative against both incumbents. If the top posterior then
recruits at both strengths, entry is flat in incumbent strength.

The next theoretical step is the seller's choice of terms with the whole pricing and trading continuation retained.
The empirical test is whether further preparation in visible sale processes responds to incumbent strength when the
price precedes the decision. A second test is whether purchase impacts carry the entry component that
Section~\ref{sec:empirical} predicts.
